% ===========================================================================
% PART 1 -- the object, and the chain, layer by layer.
% ===========================================================================
\section{The object}
\label{sec:object}

\subsection{What a normative disposition is, formally}

\begin{definition}[Disposition]
Let $\resp$ be the set of all possible responses to a prompt -- unbounded. A
person $i$'s \emph{normative disposition} is a preference structure $\succeq_i$
on $\resp$. It is not a number, not a sentence, and not a ranking: it is a
relation on an infinite set.
\end{definition}

\begin{definition}[Elicitation]
An \emph{elicitation} observes $\succeq_i$ only on a finite \emph{probe set}
$X = \{x_1,\dots,x_k\} \subset \resp$. In the object studied here, $k=4$.
\end{definition}

\subsection{The chain}

\begin{equation}
\succeq_i \;\to\; \Nn_i \;\to\; \Agg \;\to\; G \;\to\; C \;\to\; J \;\to\; D \;\to\; Y
\label{eq:chain}
\end{equation}
\noindent
$\Nn_i$ the elicited individual norm; $\Agg$ the aggregation rule; $G=\Agg(\Nn_1,\dots,\Nn_n)$
the collective norm; $C$ the compiled standard; $J$ the executor that reads it;
$D$ the decision it induces; $Y$ the behaviour of a system trained under it.

\begin{remark}[Three kinds of loss, which must never be merged]
\label{rem:threelosses}
\begin{enumerate}
\item \textbf{Mathematically forced}: no pipeline can avoid it (\cref{thm:probe}).
\item \textbf{Authorised social choice}: aggregation necessarily discards
      individual information (\cref{thm:garble}). This is a \emph{defect} only if
      the rule $\Agg$ was never declared.
\item \textbf{Machine distortion}: an unlicensed change.
\end{enumerate}
The released object declares no $\Agg$. \emph{Therefore categories 2 and 3 are
not separable at this site.} That is a missing column, not a hard measurement
problem, and it is the single most consequential fact about the release.
\end{remark}

\subsection{The tensors, exactly}

\begin{definition}[The released quantities]
\label{def:tensors}
For prompt $p$, criterion index $c$, response letter $r \in \{A,B,C,D\}$,
annotator $a$:
\begin{align*}
\sat[p,c,r] &\in (0,1) && \text{satisfaction of criterion } c \text{ by response } r;\\
\wt[p,c,a] &\in [-10,10]\setminus\{0\} && \text{the weight } a \text{ gave } c;\\
V[p,a] &\subseteq \{A,B,C,D\} && \text{the responses } a \text{ called unacceptable};\\
\Pi[p,a] &\in \text{rankings of } \{A,B,C,D\} && \text{the ranking } a \text{ gave.}
\end{align*}
\end{definition}

Four facts about \cref{def:tensors} are load-bearing and are stated here so that
no later section can assume otherwise.

\begin{enumerate}
\item $\sat$ is \emph{not} in the release. It is reconstructed by running a
      language model as a judge; the value is a sigmoid of a logit difference and
      has never been calibrated against anything. Every number downstream is a
      number about that judge as much as about the norm (\cref{thm:atten},
      §\ref{sec:instrument}).
\item $c$ is \emph{ragged}: between $4$ and $39$ criteria per prompt, median
      $15$. Any statistic that averages over criteria therefore weights prompts
      unequally unless it is explicitly re-weighted.
\item $\wt$ is \emph{sparse}: $102{,}147$ observed $(c,a)$ pairs. Exactly one
      rating in the corpus equals $0$, so the scale is used as strictly two-sided.
\item \textbf{$V$ and $\wt$ do not share an index.} $V$ is indexed by
      $(\text{prompt},\text{annotator})$; $\wt$ by
      $(\text{prompt},\text{criterion},\text{annotator})$. There is no column
      linking a veto to a criterion. §\ref{sec:type} shows what this forbids.
\end{enumerate}

\begin{remark}[The index itself is reconstructed]
The rubric file carries no prompt identifier. The mapping from criteria to
prompts is recovered by matching conversation text, with a fuzzy fallback:
$966$ exact, $2$ fuzzy, and $18$ of $986$ records \textbf{never matched at all}.
So $1.8\%$ of the rubric corpus is invisible to every number computed from this
release, including every number in this paper.
\end{remark}

\subsection{What a claim about this object can mean}

\begin{definition}[Induced score]
\label{def:score}
$\displaystyle Y_p(r) \;:=\; \sum_{c} \bar w_{p,c}\,\sat[p,c,r]$, where
$\bar w_{p,c}$ is the mean over annotators of $\wt[p,c,a]$.
\end{definition}

\begin{definition}[Individual induced score]
\label{def:ni}
$\displaystyle \Nn_{i,p}(r) \;:=\; \sum_{c \,:\, i \text{ rated } c} \wt[p,c,i]\,\sat[p,c,r]$.
\end{definition}

\begin{remark}[Why $\Nn_i$ must be centred and normalised, and why that is forced]
\label{rem:normalise}
Only contrasts among the four responses decide anything, so adding a constant to
$\Nn_i$ changes nothing; the mean over $r$ is therefore unidentified and must be
removed. Furthermore, $\|\Nn_i\|$ grows mechanically with the number of criteria
$i$ rated: measured
$\Corr(\log\|\Nn_i\|,\ \log \#\{\text{criteria rated}\}) = +0.842$. Any statistic
on unnormalised $\Nn_i$ therefore measures \emph{effort}, not values. Centring and
normalising are not stylistic; without them the quantity is a different quantity.
\end{remark}

% ===========================================================================
\section{The elicitation layer}
\label{sec:elicit}

\subsection{The probe bound}

\begin{theorem}[Probe bound]
\label{thm:probe}
Let $u : X \to \R$ be a utility on $k$ probe responses. If decisions depend only
on differences of $u$, the space of decision-distinct utilities has dimension
$k-1$. If in addition only the ordering matters (no interpersonal scale), the
space is $S^{k-2}$, of dimension $k-2$.
\end{theorem}
\begin{proof}
\emph{Step 1.} The set of utilities is $\R^k$.

\emph{Step 2 (quotient by constants).} $u$ and $u + c\mathbf{1}$ induce identical
differences $u(x_a)-u(x_b)$ for every pair, hence identical decisions. The
subspace of constants $\R\mathbf{1}$ has dimension $1$, so the quotient
$\R^k/\R\mathbf{1}$ has dimension $k-1$. Concretely, the centring map
$C = I - \tfrac1k\mathbf{1}\mathbf{1}^\top$ is a projection with
$\operatorname{rank}(C) = k-1$; for $k=4$ this rank is verified numerically as
$3$.

\emph{Step 3 (quotient by scale).} If $u$ and $\lambda u$ for $\lambda>0$ are
also equivalent, we quotient the $(k-1)$-dimensional centred space by the
positive rays, giving the sphere $S^{k-2}$, of dimension $k-2$.

\emph{Step 4.} For $k=4$: dimension $3$ after centring, and $S^2$ -- two free
parameters -- after scale.
\end{proof}

\mesotakeaway{However many criteria a person writes and however many weights they
attach, the elicitation returns \textbf{a point on a 2-sphere}: two free numbers
per person per prompt. This is arithmetic, not a finding, and it is the ceiling
on everything downstream.}

\begin{corollary}[Three consequences]
\label{cor:probe3}
\begin{enumerate}
\item The numerator of ``how much normative information did the pipeline
      destroy'' is bounded by two parameters. Most of what a criterion says never
      entered the measurable object at all.
\item Any measure of loss based on \emph{rank} or \emph{dimension} must read a
      constant. An earlier round of this work built such a measure; it returned
      the value $3.00$ on every input including its own null. That was not a bug
      -- it was \cref{thm:probe}.
\item Loss must be measured as \emph{angle} and \emph{influence}. You cannot lose
      a dimension you never had; you can only be rotated away from what people
      wanted.
\end{enumerate}
\end{corollary}

\subsection{Is the weight scale ordinal or interval? A test, and an answer}
\label{sec:cardinal}

Nothing in the elicitation establishes that the distance from $+8$ to $+10$ means
the same as from $+2$ to $+4$. If the scale is merely \emph{ordinal} -- only the
ordering of criteria within an annotator is meaningful -- then any strictly
increasing transformation of the weights must leave every decision unchanged.

\begin{protocol}[Monotone gauge test]
Apply $\varphi$ with $\varphi$ strictly increasing and odd, recompute
$Y_p$ from \cref{def:score} with $\varphi(\bar w)$ in place of $\bar w$, and
record whether $\argmax_r Y_p(r)$ moves.
\end{protocol}

\begin{table}[H]\centering\small
\begin{tabular}{lr}
\toprule
transform $\varphi(w)$ & winner changes \\
\midrule
identity (control) & $0.0\%$ \\
$\operatorname{sign}(w)|w|^{1/2}$ & $4.4\%$ \\
$\operatorname{sign}(w)\log(1+|w|)$ & $4.6\%$ \\
$\operatorname{sign}(w)|w|^{2}$ & $7.3\%$ \\
$\operatorname{sign}(w)$ alone & $\mathbf{11.7\%}$ \\
\bottomrule
\end{tabular}
\caption{Every transform preserves each annotator's ordering of criteria exactly.
If the scale were ordinal, all rows would read $0.0\%$. $n = 968$ prompts.}
\label{tab:cardinal}
\end{table}

\mesoboundary{The scale is \emph{used} as an interval scale. $7.3\%$ of decisions
change under a transform that preserves every stated preference ordering, and
$11.7\%$ change when magnitude is discarded entirely. Every result in this
literature that sums weighted criteria depends on an interval assumption the
elicitation never established. The last row is the honest headline: roughly one
decision in eight is carried by magnitude rather than direction.}

\subsection{What the elicitation captured: consensus, not the person}

\begin{table}[H]\centering\small
\begin{tabular}{lr}
\toprule
quantity ($14{,}764$ person--prompt pairs) & value \\
\midrule
$\cos(\Nn_i,\ \text{person } i\text{'s own ranking})$ & $+0.3928$\\
$\cos(\Nn_i,\ \text{a \emph{different} person's ranking})$ & $+0.3097$\\
difference & $+0.0831$ (se $0.0067$, $z=+12.3$)\\
\bottomrule
\end{tabular}
\caption{Both vectors centred and normalised per \cref{rem:normalise}.}
\label{tab:selfother}
\end{table}

Two readings, and the second is larger than the first.

\begin{itemize}
\item \textbf{Relative.} The entire individual signal is $0.083$ of cosine. A
      person's written criteria predict a \emph{stranger's} ranking almost as
      well as their own.
\item \textbf{Absolute.} $0.393$ means a person's own stated criteria recover
      only about $39\%$ of the direction of their own choice.
\end{itemize}

\mesotakeaway{Before the pipeline does anything, roughly $61\%$ of the relation
between what a person \emph{writes} and what they \emph{choose} is already
absent. No downstream step can destroy individual normative content that the
elicitation never captured; ``how much of a person survives compilation'' is
bounded above by how much of them was ever written down.}

\subsection{The normative payload actually collected}

\begin{table}[H]\centering\small
\begin{tabular}{lrr}
\toprule
field, by lexical marker & criteria & share \\
\midrule
exception / carve-out (\emph{unless, except, other than}) & $123$ & $0.8\%$\\
conditional scope (\emph{if, when, where, in cases of}) & $1{,}724$ & $11.3\%$\\
priority / conflict (\emph{before, over, takes precedence}) & $440$ & $2.9\%$\\
uncertainty (\emph{may, might, unclear, depends}) & $478$ & $3.1\%$\\
non-compensable (\emph{never, must not, regardless}) & $121$ & $0.8\%$\\
\midrule
\textbf{none of the five} & $\mathbf{12{,}547}$ & $\mathbf{82.3\%}$\\
\bottomrule
\end{tabular}
\caption{All $15{,}248$ criteria. Median length $94$ characters.}
\label{tab:payload}
\end{table}

A minimal normative payload has been proposed with thirteen fields: content,
scope, polarity, exceptions, priority, type, compensability, authority, judgment
type, uncertainty, witnesses, constituency, meta-rules. This release populates
\textbf{three}: content, polarity (the sign of a weight), magnitude.

\begin{corollary}[A question that is empty here]
``Did the compilation destroy the exception clause?'' cannot be studied at this
site: $99.2\%$ of criteria contain no exception. A pipeline cannot destroy
structure the elicitation never collected. This is a statement about the release,
not about pipelines.
\end{corollary}

% ===========================================================================
\section{The aggregation layer}
\label{sec:agg}

\subsection{Social choice, from zero}

Given $n$ people each with a ranking of $k$ alternatives, a \emph{social choice
rule} outputs one alternative. Five rules appear below; each is defined
completely.

\begin{definition}[Rules]
\label{def:rules}
Let $\pi_i(x)$ be the Borda points person $i$ gives $x$ (i.e. $k-1$ for their
first choice down to $0$ for last, ties sharing the average).
\begin{itemize}
\item \textbf{Plurality}: pick the $x$ that is ranked first by the most people.
\item \textbf{Borda}: pick $\argmax_x \frac1n\sum_i \pi_i(x)$.
\item \textbf{Median}: pick $\argmax_x \operatorname{median}_i \pi_i(x)$.
\item \textbf{Maximin}: pick $\argmax_x \min_i \pi_i(x)$ -- the alternative whose
      \emph{worst} treatment by any individual is least bad.
\item \textbf{Condorcet winner}: the $x$ that beats every other alternative in
      pairwise majority comparison. It need not exist.
\end{itemize}
\end{definition}

\begin{definition}[Condorcet cycle]
A majority cycle is a triple with $x$ beating $y$, $y$ beating $z$, and $z$
beating $x$. Majority preference is then not transitive and no Condorcet winner
exists.
\end{definition}

\subsection{Measured social choice structure}

\begin{table}[H]\centering\small
\begin{tabular}{lr}
\toprule
quantity ($968$ prompts, $\ge 3$ rankings each) & value \\
\midrule
a Condorcet winner exists & $72.3\%$\\
a majority \emph{cycle} occurs & $\mathbf{0.0\%}$\\
neither (majority ties) & $27.7\%$\\
\bottomrule
\end{tabular}
\caption{The classic voting pathology does not occur in this corpus.}
\label{tab:condorcet}
\end{table}

\begin{table}[H]\centering\small
\begin{tabular}{lr}
\toprule
agreement with the Condorcet winner, where one exists ($n=700$) & \\
\midrule
Borda & $98.4\%$\\
Plurality & $98.3\%$\\
Median & $97.4\%$\\
\textbf{the weighted-criteria score of \cref{def:score}} & $\mathbf{68.4\%}$\\
Maximin & $55.9\%$\\
\bottomrule
\end{tabular}
\label{tab:agree}
\end{table}

\mesotakeaway{Every standard \emph{ranking-based} rule agrees with the Condorcet
winner to $97$--$98\%$. The weighted-criteria score agrees only $68.4\%$ of the
time. The disagreement is therefore \textbf{not} between aggregation rules; it is
between \emph{scoring people's criteria} and \emph{counting people's rankings} --
two elicitations from the same people whose outputs differ on nearly a third of
prompts.}

\subsection{Arrow's independence condition, measured}

\begin{definition}[IIA]
A rule satisfies \emph{independence of irrelevant alternatives} if the social
ordering of $x$ and $y$ depends only on individual orderings of $x$ and $y$, not
on any third alternative.
\end{definition}

Measured for Borda: over $11{,}616$ tests (each prompt $\times$ each dropped
response $\times$ each remaining pair), the relative order of two responses
reverses when an irrelevant third is removed in $1{,}553$ cases:
$\mathbf{13.4\%}$.

\begin{corollary}[The candidate set is part of the decision]
\label{cor:iia}
$13.4\%$ of pairwise relations flip when the option set changes. Any normative
content measured on one set of four responses is, to that extent, a property of
the set. This is the quantitative content of ``candidate-set overfitting''.
\end{corollary}

\subsection{Blackwell: what aggregation necessarily costs}

\begin{definition}[Garbling]
$G$ is a \emph{garbling} of $Z$ if there is a Markov kernel $K$ with
$G \sim K(\cdot \mid Z)$: $G$ is generated from $Z$ by a channel that adds no new
information about the state.
\end{definition}

\begin{lemma}[Aggregation is a garbling]
\label{lem:garble}
If $G = \Agg(Z)$ is a deterministic function of $Z$, then $G$ is a garbling of
$Z$, via the degenerate kernel $K(\cdot\mid z) = \delta_{\Agg(z)}$.
\end{lemma}

\begin{theorem}[Blackwell, invoked]
\label{thm:garble}
If $G$ is a garbling of $Z$ then for \emph{every} decision problem $d$ with any
loss function, the attainable Bayes risk satisfies $R_G(d) \ge R_Z(d)$.
\end{theorem}

\begin{corollary}
Aggregation never adds information and weakly loses on every decision problem.
The question is never \emph{whether}, only \emph{how much} -- and ``how much'' is
meaningless until a decision family is named.
\end{corollary}

\begin{definition}[Deficiency]
$\displaystyle \delta_{\mathcal{D}}(Z,G) := \sup_{d \in \mathcal{D}}\big[R_G(d)-R_Z(d)\big]$.
\end{definition}

\step{Computing it here.} Take $\mathcal{D}$ to be the single problem ``choose one
of the four responses''. Normalise each person's induced score to $[0,1]$ by
min--max within that person and prompt, so that scores are interpersonally
comparable by construction rather than by assumption. Then:
\begin{align*}
\text{oracle (each person receives their own top response)} &= 1.0000 \\
\text{the single collective choice} &= 0.8702 \\
\textbf{deficiency} &= \mathbf{0.1298}
\end{align*}
with quartiles $0.0445 / 0.1076 / 0.1924$ across prompts.

\mesotakeaway{Thirteen points of normalised utility is what \emph{disagreement
itself} costs. It is a social choice, not a pipeline defect, and no compilation
recovers it. It also fixes the denominator: ``what fraction was preserved''
should be divided by $0.8702$, not by $1$.}

\subsection{Sen's liberal paradox, instantiated}

\begin{table}[H]\centering\small
\begin{tabular}{lr}
\toprule
prompts carrying a partial veto & $161$\\
the weighted-criteria winner \emph{is} a vetoed response & $42 = \mathbf{26.1\%}$\\
\bottomrule
\end{tabular}
\label{tab:sen}
\end{table}

One time in four, a rule that maximises aggregate satisfaction selects an option
someone declared unacceptable. This is not a defect of the rule; it is what a
utilitarian aggregator does. It is the empirical half of §\ref{sec:type}.

\subsection{Which rule is the larger lever, corrected}

An earlier version of this analysis applied ``median'' and ``maximin'' across
\emph{criteria}. A social choice rule aggregates across \emph{people}; the two
coincide only if each criterion belongs to exactly one person, which is false
here. Corrected:

\begin{table}[H]\centering\small
\begin{tabular}{lrr}
\toprule
rule, versus the mean over people & winner changes & (wrong index: over criteria)\\
\midrule
maximin & $49.4\%$ & $63.9\%$\\
median & $14.8\%$ & $34.3\%$\\
Borda & $9.6\%$ & --\\
trimmed & $\mathbf{3.6\%}$ & --\\
\midrule
\emph{any intervention on a single criterion} & \emph{$2.1$--$11.6\%$} &\\
\bottomrule
\end{tabular}
\label{tab:aggrule}
\end{table}

\begin{remark}[And maximin needs a control]
A minimum over a set is unstable by construction. Control: same people, same
per-person mean and standard deviation, response pattern destroyed. That sham
disagrees with the mean rule on $61.4\%$ of prompts -- \emph{more} than the real
data's $49.4\%$. So maximin's high rate is the instability of the $\min$
operator, and the actual structure of human disagreement \emph{reduces} it by
$12$ points. The claim ``minority-sensitive rules change half the outcomes'' does
not survive its own control. What survives is that \emph{trimmed} aggregation,
at $3.6\%$, moves less than several criterion-level interventions: trimming the
extremes barely matters, so the disagreement that matters lives in the tails.
\end{remark}

% ===========================================================================
\section{The type layer: why a veto is not a large negative weight}
\label{sec:type}

\subsection{The index mismatch}

From \cref{def:tensors}: $V$ is indexed $(p,a)$ and $\wt$ is indexed $(p,c,a)$.
An operator ``represent this veto as a weight'' must choose \emph{which criterion}
carries it. No such column exists.

\begin{proposition}
``Coerce a veto into a criterion weight'' is not a function of the released data.
\end{proposition}

This is a schema fact, not a measurement difficulty, and no amount of data at
this site repairs it.

\subsection{The well-posed question, and its answer}

\begin{definition}[Two rules]
\begin{align}
\text{constraint:}&\quad a^\star \;=\; \argmax_{a \notin V} Y(a),\\
\text{penalty:}&\quad a^\star_M \;=\; \argmax_{a} \big[Y(a) - M\cdot\mathbf{1}\{a\in V\}\big].
\end{align}
\end{definition}

\begin{theorem}[Exact condition for a penalty to reproduce a veto]
\label{thm:veto}
Let $\Delta := \max_{a\in V} Y(a) - \max_{a\notin V} Y(a)$. Then
$a^\star_M = a^\star$ for every realisation of the scores if and only if
$M > \Delta$.
\end{theorem}
\begin{proof}
\emph{($\Leftarrow$)} Suppose $M > \Delta$. For any $a \in V$,
$Y(a) - M \le \max_{a'\in V}Y(a') - M < \max_{a'\in V}Y(a') - \Delta
= \max_{a'\notin V}Y(a')$. So every vetoed alternative is strictly below the best
allowed one after the penalty, and the $\argmax$ lies in the allowed set, where
the penalty is zero and the ordering is unchanged. Hence $a^\star_M = a^\star$.

\emph{($\Rightarrow$)} Suppose $M \le \Delta$. Let $a_V$ attain
$\max_{a\in V}Y(a)$. Then
$Y(a_V) - M \ge Y(a_V) - \Delta = \max_{a\notin V}Y(a)$, so $a_V$ is weakly
preferred by the penalty rule to every allowed alternative, and the penalty rule
may select a vetoed response. Hence agreement fails.
\end{proof}

\begin{table}[H]\centering\small
\begin{tabular}{lr}
\toprule
measured quantity & value\\
\midrule
prompts with a partial veto & $161$\\
the veto is \emph{binding} ($\Delta > 0$) & $42 = 26.1\%$\\
required $M$: median & $3.64$\\
required $M$: 90th percentile & $9.92$\\
required $M$: \textbf{maximum} & $\mathbf{26.54}$\\
contribution of one criterion at the scale maximum & $\le 10$\\
required $M$ in units of max-weight criteria: median / max & $0.36$ / $\mathbf{2.65}$\\
\bottomrule
\end{tabular}
\label{tab:vetoM}
\end{table}

\begin{theorem}[No finite weight encodes a veto]
\label{thm:vetoinf}
$\Delta$ is a function of the \emph{candidate set}, not of the norm. Since a
vetoed response can be arbitrarily better than every allowed one on the remaining
criteria, $\sup_X \Delta(X) = \infty$. Therefore no finite $M$ satisfies
\cref{thm:veto} uniformly over candidate sets.
\end{theorem}
\begin{proof}
Fix a person, a veto set $V$, and a criterion $c$ with weight $w>0$. Construct a
candidate $x^\ast \in V$ with $\sat[\cdot,c,x^\ast] \to 1$ and all allowed
candidates with $\sat[\cdot,c,\cdot] \to 0$; then $\Delta \to w \cdot \sum$ over
however many such criteria are present, which is unbounded as the number of
agreeing criteria grows. Since \cref{thm:veto} requires $M>\Delta$ for agreement
on that set, no fixed finite $M$ suffices for all sets.
\end{proof}

\mesotakeaway{A veto is a property of the \emph{person}. The penalty $M$ that
reproduces it is a property of the person \textbf{and the four responses that
happened to be shown}. On just $968$ four-item candidate sets the required
penalty already reaches $2.65$ times the entire elicitation scale. Type is not
compressible to magnitude, and this is a theorem with its counterexample present
in the data -- it required no instrument at all.}

\subsection{Three independent routes to one conclusion}

\begin{table}[H]\centering\small
\begin{tabular}{lll}
\toprule
route & evidence & section\\
\midrule
derivation & no finite $M$ works uniformly & \cref{thm:vetoinf}\\
empirical & the utilitarian winner is vetoed $26.1\%$ of the time & \cref{tab:sen}\\
structural & the index does not exist in the schema & §\ref{sec:type}\\
\bottomrule
\end{tabular}
\end{table}

% ===========================================================================
\section{The compilation layer}
\label{sec:compile}

The released standard ships in two forms: the full weighted rubric and a compiled
``core'' whose items carry no weights, every one rewritten to positive polarity.

\begin{itemize}
\item The two select \emph{different} winners on $\mathbf{29.4\%}$ of prompts.
\item Compilation performs three operations at once -- selection, rewriting, and
      weight discarding -- and the release ships no lineage, so the three are
      \textbf{not separable}. This is the missing column that blocks the most
      questions here.
\end{itemize}

\subsection{How many bits of the rubric are decision-relevant?}

\begin{protocol}[Rate--distortion]
Quantise every criterion weight to $b$ bits uniformly across the observed range
within each prompt, recompute \cref{def:score}, and record whether the winner is
preserved.
\end{protocol}

\begin{table}[H]\centering\small
\begin{tabular}{lr}
\toprule
bits per criterion weight & decision preserved\\
\midrule
$1$ (sign only) & $\mathbf{89.8\%}$\\
$2$ & $95.0\%$\\
$3$ & $98.1\%$\\
$4$ & $98.6\%$\\
$8$ & $99.8\%$\\
\bottomrule
\end{tabular}
\label{tab:rd}
\end{table}

\begin{remark}[Two independent computations agree]
\cref{tab:rd} says one bit preserves $89.8\%$, i.e. changes $10.2\%$.
\cref{tab:cardinal} says discarding magnitude changes $11.7\%$. These were
computed by different routes and agree to about a point. The $21$ levels of the
$-10..{+}10$ scale carry roughly three bits of decision-relevant content, and the
\emph{first} bit carries nine tenths of it.
\end{remark}

\begin{corollary}[A design consequence]
The elicitation is over-precise in magnitude and empty in type
(\cref{tab:payload}). The next collection should not buy finer weights; it should
buy fields.
\end{corollary}

% ===========================================================================
\section{The instrument layer}
\label{sec:instrument}

\subsection{The judge as a noisy channel}

$\sat$ is produced by a language model. Two independent judges give two
measurements of the same underlying satisfaction, so \cref{eq:relcorr} applies
directly.

\begin{table}[H]\centering\small
\begin{tabular}{lr}
\toprule
$59{,}936$ satisfaction cells & value\\
\midrule
$\Corr(\text{judge}_1,\text{judge}_2)$ & $0.5497$\\
$\Rightarrow$ reliability of one judge (\cref{eq:relcorr}) & $0.5497$\\
$\Rightarrow$ attenuation factor $\sqrt{\mathrm{rel}}$ (\cref{thm:atten}) & $0.741$\\
\bottomrule
\end{tabular}
\label{tab:channel}
\end{table}

\begin{corollary}[A ceiling on every number in this literature]
A true correlation of $1.0$ measured through this tensor reads at most $0.741$.
This is bias, not noise: more prompts do not remove it.
\end{corollary}

\subsection{Panel reliability}

Split-half over annotators on $964$ prompts with $\ge 6$ annotators:
$r = 0.8961$; by \cref{thm:sb}, full-panel reliability
$2r/(1+r) = 0.9452$, so the panel score alone caps correlations at
$\sqrt{0.9452} = 0.9722$.

\begin{corollary}[Compounded ceiling]
A correlation measured through both the panel and the judge is capped at
$0.9722 \times 0.741 = 0.720$.
\end{corollary}

\subsection{The gauge test}
\label{sec:gauge}

\begin{definition}[Gauge transformation]
A transformation that must leave the \emph{property} unchanged. Here: change the
reader of the norm while holding the norm fixed. If the \emph{measurement} is
invariant but the property is not, the measurement is blind; if the property is
invariant but the measurement is not, the measurement is about the reader.
\end{definition}

\begin{table}[H]\centering\small
\begin{tabular}{lrr}
\toprule
& winner differs from baseline & $n$\\
\midrule
judge B (different family) & $42.4\%$ & $968$\\
judge C (different scale) & $39.7\%$ & $968$\\
\textbf{same judge, responses reordered} & $\mathbf{8.0\%}$ & $300$\\
\bottomrule
\end{tabular}
\caption{(a) The object is \emph{not} invariant.}
\label{tab:gaugea}
\end{table}

\begin{table}[H]\centering\small
\begin{tabular}{lrrr}
\toprule
Spearman $\rho$ of the operator ordering, against baseline & judge B & judge C & reordered\\
\midrule
& $+0.982$ & $+1.000$ & $+0.963$\\
\bottomrule
\end{tabular}
\caption{(b) The measurement \emph{is} invariant. Pre-registered kill at
$\rho<0.8$; it does not fire on any of the three.}
\label{tab:gaugeb}
\end{table}

\mesotakeaway{Comparative statements are about the norm. Absolute statements are
about the reader. ``Operator $X$ moves the decision more than $Y$'' survives three
instruments and a position swap. ``The standard prefers response $X$'' does not:
roughly $40\%$ of it is the reader and $8$ points are nothing but the order the
responses were printed in.}

\begin{remark}[The scale reference, with its caveat]
Reordering the page changes the winner more often than deleting a criterion does
($8.0\%$ versus $5.1\%$). The caveat is load-bearing: the reordering comparison
uses a \emph{separate} judge run and so carries judge stochasticity, whereas the
operator comparison holds the tensor fixed and cancels it. So $8.0\%$ is not a
strict noise floor for the operators; it is the correct \emph{scale reference}
for any claim of the form ``this changes the decision''.
\end{remark}
